\scanpage{46}
% sheet rotated 180 degrees
$\lambda=(\underset{\leftarrow}{\lambda},\varnothing)$ \qquad muni du foncteur de cat. 2-cofibrées $T\xrightarrow{b^T_{\unsure{\tau|?|C}}}\tau^I$

1) $T$ catégorie \unsure{2-cofibrée|1-cofibrée|$\lambda$-abstraite} sur $(\mathrm{Cat}_\lambda)$. On suppose que \deleted{pour \unsure{tte} $C$}

a) Pour tt foncteur pl. fidèle, $i:C\to C'$ dans $\mathrm{Cat}_\lambda$, $T(i):T(C)\to T(C')$ pl. fid.

\deleted{b)} et le carré

\diagram{\begin{tikzcd}[ampersand replacement=\&] T(C)\arrow[r]\arrow[d,"{b^T_C}"']\&T(C')\arrow[d,"{b^T_{C'}}"]\\ C^I\arrow[r]\&{C'}^I\end{tikzcd}}

est 2-cartésien [il \unsure{le} suffit pour $C\lhook\joinrel\xrightarrow{i'}\widehat C$]

% slanted
\margin{Attention au cas où $\widehat C$ n'est pas une $\mathcal U$-catégorie !}

b) Pour tt $C\in\mathrm{Ob}\,\mathrm{Cat}_\lambda$, le foncteur canonique

\[T(\widehat C)\longrightarrow\underline{\mathrm{Hom}}(C^\circ,\mathscr S)\qquad \mathscr S\overset{\mathrm{déf}}=T(\mathrm{Ens})\]

est une équivalence de catégories.

Alors $T\deleted{\unsure{(C)}}$ se réconstitue à partir de la donnée de $\mathscr S$ et de $(\beta=b_{\mathrm{Ens}}):\mathscr S\to(\mathrm{Ens})^I$, $(\beta=(\beta_i)_{i\in I},\ \beta_i:\mathscr S\to(\mathrm{Ens}))$ comme $T(C)=$ ss-catégorie pleine de $\underline{\mathrm{Hom}}(C^\circ,\mathscr S)$ formée des $\varphi:C^\circ\to\mathscr S$ tels que $\forall i\in I$, $\beta_i\circ\varphi$ soit représentable.

\margin{Mais si $f:C\to C'$ est un $\lambda$-foncteur, comment définir $T(f):T(C)\to T(C')$ en termes de $\mathscr S$, $(\beta_i)$ ?}

\deleted{\unsure{Il suffit}} La ss-catégorie pleine $R$ de $\underline{\mathrm{Hom}}(\mathscr S,(\mathrm{Ens}))$ formée des foncteurs $\psi\deleted{\unsure{g}}$ tels que pour tt $C\in\mathrm{Ob}\,\mathrm{Cat}_\lambda$ \deleted{\unsure{$\xi$}} et $\xi\in T(C)$, le composé $\psi\circ\xi:C^\circ\longrightarrow\mathscr S\xrightarrow{\psi}(\mathrm{Ens})$ soit représentable, est stable par $\varprojlim$ de type $\in\underset{\leftarrow}{\lambda}$, \deleted{A quelle condition} et pour tout $C\in\mathrm{Ob}\,\mathrm{Cat}_\lambda$, on a un accouplement

\[T(C)\times R\longrightarrow C\qquad\substack{\text{commutant aux }\varprojlim\text{ de}\\\text{type }\underset{\leftarrow}{\lambda}\text{ sur }R}\]

d'où $T(C)\to\operatorname{Hom}_{\lambda}(R,C)$ (fonct. de cat. 2-cof. sur $\mathrm{Cat}_\lambda$)

À quelle condition cela est-il une équivalence de catégories ?

Il est néc que a) $R$ soit \inserted{formée de fonct.} repr. i.e. $R\simeq\Sigma^\circ$, où $\Sigma$ est une sous-catégorie pleine de $\mathscr S$ [dans laquelle les $\varinjlim$ (\unsure{inductives|coinductives|inductibles}) de type $\lambda$ existent dans $\mathscr S$, et appartiennent à $\Sigma$]. Donc la donnée de $R$ équivaut à la donnée de $\Sigma\subset\mathscr S$

b) $\Sigma$ engendre $\mathscr S$ $[\,\mathscr S\to\underline{\mathrm{Hom}}(\Sigma^\circ,(\mathrm{Ens}))=\widehat\Sigma\ \text{pl. fidèle}\,]$
% archive sheet number: 45
